\cleardoublepage
\thispagestyle{plain}
\begingroup
\small
\setstretch{1.08}
\setlength{\parindent}{0pt}
\setlength{\parskip}{0.75em}
\phantomsection\addcontentsline{toc}{chapter}{How to Cite \progname{}}
\begin{center}
{\LARGE\bfseries How to Cite Diabat}\par
\end{center}
\vspace{0.6cm}
If you use \progname{} in published work, please cite the software release article. The final journal information and DOI will be added to this manual after formal publication.

\begin{quote}
Y.-C. Wang, Y. Qin, Y. Zhao, and W. Z. Liang.\\
\textit{Diabat 2.0: A Comprehensive Diabatization Toolkit for Multichromophoric Systems.}\\
\textit{Journal of Chemical Theory and Computation}, [year], [volume], [pages].\\
DOI: [DOI to be updated].
\end{quote}

If your calculations use the fragment particle-hole densities (FPHD) method, please also cite the two FPHD method papers.

\begin{quote}
Y.-C. Wang, S. Feng, W. Z. Liang, and Y. Zhao.\\
\textit{Electronic Couplings for Photoinduced Charge Transfer and Excitation Energy Transfer Based on Fragment Particle--Hole Densities.}\\
\textit{J. Phys. Chem. Lett.} \textbf{2021}, \textit{12}, 1032--1039.\\
\href{https://doi.org/10.1021/acs.jpclett.0c03514}{doi:10.1021/acs.jpclett.0c03514}

\vspace{0.45em}
Y.-C. Wang, S. Feng, Y. Kong, X. Huang, W. Z. Liang, and Y. Zhao.\\
\textit{Electronic Couplings for Singlet Fission Processes Based on the Fragment Particle--Hole Densities.}\\
\textit{J. Chem. Theory Comput.} \textbf{2023}, \textit{19}, 3900--3914.\\
\href{https://doi.org/10.1021/acs.jctc.3c00243}{doi:10.1021/acs.jctc.3c00243}
\end{quote}

References for the other diabatization methods are given in the corresponding method sections.
\par\endgroup\clearpage
